\section{Exact exceptional pin sets and the boundary}\label{sec:pins}
\subsection{Off the horizontal support}
\begin{proposition}\label{prop:offH}
Under the superlacunary assumptions, or the fixed-ratio assumptions with
$aK>1$, every $p\in P\setminus\mathcal B$ has a bounded continuous
original density. The density is jointly continuous locally on
$(P\setminus\mathcal B)\times\R$, and
\begin{equation}\label{eq:off-H}
 \norm{g_p}_\infty\leq\frac{C}{\dist(p_1,H)}.
\end{equation}
\end{proposition}
\begin{proof}
Since $H$ is compact, $d=\dist(p_1,H)>0$. Set $\varepsilon=d/2$.
On a neighborhood of $p$ with $\dist(p'_1,H)>\varepsilon$, the mask
$M_\varepsilon(p',\cdot)$ is identically one on the entire source
support. The jointly continuous positive-filter density of
Proposition~\ref{prop:fixed-filter} is therefore the original density
throughout that neighborhood, with bound $2C_*/d$.
\end{proof}
Every fixed off-$H$ pin has a positive gap, so no third membership class
arises from a rate of approximation to $H$. Bounds may diverge as pins
approach $H$, which is a quantitative issue distinct from membership.

\subsection{Dimensions of the full horizontal-support pin set}
\begin{proposition}\label{prop:B-dimensions}
In the superlacunary regime,
\begin{equation}\label{eq:B-super-dim}
 \dimH H=\alpha-1,\quad\dimP H=1,\qquad
 \dimH\mathcal B=\alpha,\quad\dimP\mathcal B=2.
\end{equation}
In the fixed-ratio regime, put
\begin{align}
 s_K&=\frac{b+aK}{1+K/2}
     =\frac{\alpha+(\alpha-1)K}{K+2},\label{eq:sK}\\
 t_K&=\frac{b+aK}{1+aK}
     =\frac{\alpha+(\alpha-1)K}{2+(\alpha-1)K}.\label{eq:tK}
\end{align}
Then
\begin{equation}\label{eq:B-fixed-dim}
 \dimH H=s_K,\quad\dimP H=t_K,\qquad
 \dimH\mathcal B=1+s_K,\quad\dimP\mathcal B=1+t_K.
\end{equation}
\end{proposition}
\begin{proof}
In the superlacunary case,~\eqref{eq:horizontal-frostman} gives the
Hausdorff lower bound for $H$. At depth $n+L/2$, the horizontal free
count is $bn+aL$, whose ratio to depth tends to $2a=\alpha-1$.
The corresponding covers give the upper bound. At depth $n+aL$ the
ratio tends to one, and every fixed cylinder has upper box dimension one.
The coding-cylinder and Baire argument of Section~\ref{sec:construction}
then gives packing dimension one.

In the fixed-ratio case the ratio
$[bn+f_X(l)]/(n+l)$ increases up to $l=aKn$, decreases up to $l=Kn/2$,
and increases thereafter. Its trough is $s_K$ and its peak is $t_K$;
indeed $s_K\leq b\leq t_K$, so no stage endpoint has been omitted.
The finitely many initial depths alter only a constant. The horizontal
law is $s_K$-Frostman, and covers at the middle-stage trough give the
matching Hausdorff upper bound. Every fixed cylinder has upper box
dimension $t_K$, and Baire gives the packing lower bound.

For $\mathcal B=H\times[0,\eta]$, multiply the appropriate horizontal
Frostman law by Lebesgue probability on the interval to obtain the
Hausdorff lower bounds. Multiply the trough covers by interval grids to
obtain the upper bounds. For packing dimension, cylinder products with
nondegenerate subintervals have upper box dimensions larger by exactly one;
the same countable-cover/Baire argument applies. This proves the product
formulas directly rather than relying on an equality for arbitrary
fractal products.
\end{proof}
Useful identities for this family are
\begin{equation}\label{eq:dimension-identities}
 D_K=(K+2)(1/2-s_K),\qquad
 q_K=\frac{1-s_K}{1-2s_K}\quad(D_K>0).
\end{equation}
The superlacunary formula has the same expression in $s=\alpha-1$,
although the endpoint behavior differs.
At $D_K=0$,
\begin{equation}\label{eq:boundary-dimensions}
 s_K=1/2,\quad t_K=\frac{K+1}{K+2},\qquad
 \dimH\mathcal B=3/2,\quad\dimP\mathcal B=2-\frac1{K+2}.
\end{equation}
For $K=9$, $\alpha=29/20$, these dimensions are $3/2$ and $21/11$.
Under the standing requirement $aK>2$, an integer boundary example must
have $K\geq9$.

\subsection{No continuous version at the fixed-ratio boundary}
\begin{proposition}\label{prop:boundary-discontinuity}
If $aK>2$ and $D_K=0$, then every $p\in\mathcal B$ has no continuous
original distance density on $\R$.
\end{proposition}
\begin{proof}
The minimum vertical source coordinate is $r_0=1-\eta$. For
$p_1\in H$, the point $(p_1,r_0)$ lies in the source support, so
\[
 r_{\min}(p)=r_0-p_2>0
\]
is the minimum source distance. Let $\delta_j=2^{-(n+L)}$.
Combine the horizontal track in~\eqref{eq:track-mass} with the lowest
occupied vertical interval of length $\delta_j$, which has
$\zeta_+$-mass $2^{k-b(n+L)}$. The product subset has mass
\begin{equation}\label{eq:boundary-mass}
 2^{2k-2bn-(a+b)L}.
\end{equation}
Its vertical distances are within $\delta_j$ of $r_{\min}(p)$ and its
horizontal quadratic error is at most $C2^{-2n-L}\leq C\delta_j$.
The actual distance image lies in
$[r_{\min}(p),r_{\min}(p)+C_0\delta_j]$. Dividing
\eqref{eq:boundary-mass} by $\delta_j$ gives
\begin{equation}\label{eq:boundary-ratio}
 2^{2k}2^{(1-\alpha)n+(3/2-\alpha)L}=2^{2k}2^{D_Kn}.
\end{equation}
At $D_K=0$ this is a fixed positive constant, so
\begin{equation}\label{eq:boundary-local}
 \nu_p([r_{\min}(p),r_{\min}(p)+C_0\delta_j])\geq c_0\delta_j.
\end{equation}
A continuous density for $\nu_p$ must vanish to the left of its support
and hence at $r_{\min}(p)$. Its supremum on the shrinking displayed
interval would tend to zero, giving mass $o(\delta_j)$, contrary to
\eqref{eq:boundary-local}.
\end{proof}
This proof does not exclude an essentially bounded density with a jump at
the support endpoint. It settles continuity, not $L^\infty$ membership.

\subsection{The exact classification}
Define
\begin{align*}
 \Bad_q&=\{p\in P:\nu_p\text{ has no }L^q\text{ density}\},\\
 \Bad_{C^0}&=\{p\in P:\nu_p\text{ has no continuous density on }\R\}.
\end{align*}
All these laws are absolutely continuous. Compact support implies that
every continuous-density version is bounded. The upper bounds,
\eqref{eq:track-lower},~\eqref{eq:fixed-lower},
Proposition~\ref{prop:offH}, and Proposition~\ref{prop:boundary-discontinuity}
therefore give the following exact statements.
\begin{enumerate}[leftmargin=*]
\item Superlacunary, $1<\alpha<3/2$:
\[
 \Bad_q=\varnothing\ (1\leq q\leq q_*),\qquad
 \Bad_q=\mathcal B\ (q_*<q\leq\infty),\qquad
 \Bad_{C^0}=\mathcal B.
\]
\item Superlacunary, $3/2\leq\alpha<2$: all $\Bad_q$,
$1\leq q\leq\infty$, and $\Bad_{C^0}$ are empty.
\item Fixed ratio, $aK>2$ and $D_K>0$:
\[
 \Bad_q=\varnothing\ (1\leq q<q_K),\qquad
 \Bad_q=\mathcal B\ (q_K\leq q\leq\infty),\qquad
 \Bad_{C^0}=\mathcal B.
\]
\item Fixed ratio, $aK>2$ and $D_K<0$: all these exceptional sets
are empty.
\item Fixed ratio, $aK>2$ and $D_K=0$:
\[
 \Bad_q=\varnothing\ (1\leq q<\infty),\qquad
 \Bad_{C^0}=\mathcal B,\qquad \Bad_\infty\subseteq\mathcal B.
\]
The exact set $\Bad_\infty$ in the last line is not determined here.
\end{enumerate}
The nonempty exact exceptional sets have the dimensions in
Proposition~\ref{prop:B-dimensions}. They are exceptions to stronger
\emph{density regularity}, not pins with zero-measure or interval-free
distance sets. Every pin in $P$ still has the interval conclusion.

\subsection{An upper bound on all boundary norms}
\begin{theorem}\label{thm:log-upper}
Assume $aK>2$ and $D_K=0$. Then for a fixed $C<\infty$,
\begin{equation}\label{eq:log-upper}
 \sup_{p\in P}\norm{g_p}_q\leq C\log(e+q)
 \qquad(1\leq q<\infty).
\end{equation}
There are $c_1,c_2,T_0>0$ such that
\begin{equation}\label{eq:double-tail}
 \sup_{p\in P}|\{r:g_p(r)>T\}|
 \leq \exp[-c_1\exp(c_2T)]\qquad(T\geq T_0).
\end{equation}
For some $c_3>0$,
\begin{equation}\label{eq:double-integrability}
 \sup_{p\in P}\int_I\exp(\exp(c_3g_p(r)))\,dr<\infty.
\end{equation}
Only upper bounds are asserted; no matching lower growth in $q$ or tail
asymptotic is proved.
\end{theorem}
\begin{proof}
The initial density and the summable outer increments
\eqref{eq:fixed-outer} have $L^q$ bounds independent of $q\geq1$, using
the fixed finite output interval. The constant in~\eqref{eq:MH} can also
be chosen uniformly in $q$: the coarea factor is
$(D_0/v_0)^{1-1/q}$ and fixed prefix factors are bounded by their
$q=\infty$ values. At $D_K=0$, the inner increment is at most
$C2^{-c(K+1)n/q}$. Thus
\begin{equation}\label{eq:lacunary-moment-series}
 \norm{g_p}_q\leq C+C\sum_{j\geq2}
 2^{-c(K+1)n_1(K+1)^{j-2}/q}.
\end{equation}
There are $O(1+\log q)$ terms before the exponent is of order one;
after that point the remaining sum is bounded by a fixed convergent
series of the form $\sum_{l\geq0}\exp[-A(K+1)^l]$.
This proves~\eqref{eq:log-upper}.

For large $T$ choose the real moment order $q=\exp[T/(4C)]$,
increasing $C$ and the lower threshold if necessary. Then
$C\log(e+q)\leq T/2$ and Markov gives
\[
 |\{g_p>T\}|\leq(\norm{g_p}_q/T)^q\leq2^{-q},
\]
which proves~\eqref{eq:double-tail}. Choose $0<c_3<c_2$ sufficiently
small. Layer-cake integration against the derivative of
$\exp(\exp(c_3T))$ gives~\eqref{eq:double-integrability}, since the
negative term $-c_1\exp(c_2T)$ dominates $\exp(c_3T)$ and the remaining
linear terms. The integral is over $I$, not all of $\R$, so its nonzero
constant baseline causes no divergence.
\end{proof}
The double-exponential upper tail neither proves nor disproves essential
boundedness. This completes all assertions of Theorem~\ref{thm:main-fixed}
and the full pin classification, with the stated boundary question intact.
